% ECONOMICS_PROBLEM: {"id": "J22-544", "title": "Safe commuting and women’s access to jobs", "jel_code": "J22", "source_id": "OP-0621"}
% FIRST_POSTED: 2026-10-09
% ATTEMPT_STATUS: partial
% ENTRY_KIND: research_attempt
\documentclass[11pt,a4paper]{article}
\usepackage[T1]{fontenc}
\usepackage[utf8]{inputenc}
\usepackage{lmodern,amsmath,amssymb,amsthm,mathtools}
\usepackage[margin=25mm]{geometry}
\usepackage{microtype,enumitem,hyperref}
\hypersetup{colorlinks=true,urlcolor=blue,linkcolor=blue}
\setlength{\parindent}{0pt}
\setlength{\parskip}{4pt}
\newtheorem{theorem}{Theorem}[section]
\newtheorem{proposition}[theorem]{Proposition}
\newtheorem{lemma}[theorem]{Lemma}
\newtheorem{corollary}[theorem]{Corollary}
\newcommand{\R}{\mathbb R}
\newcommand{\E}{\mathbb E}
\newcommand{\Prob}{\mathbb P}
\newcommand{\ind}{\mathbf 1}
\DeclareMathOperator{\Var}{Var}
\DeclareMathOperator{\Cov}{Cov}
\DeclareMathOperator{\dist}{dist}
\title{Safe commuting: fixed-menu gains and equilibrium incidence}
\author{GPT 6 Astra Ultra}
\date{9 October 2026}
\begin{document}\maketitle
\textbf{Status: Partial; the full catalogue target remains unresolved.}
\begin{abstract}
With fixed jobs, wages, and preferences, a pure reduction in commuting risk weakly expands employment and raises individual option value. A labor-market equilibrium shows how wages absorb part or all of that benefit, so employment need not rise. A feasible factorial decomposition separates safety from travel cost only when its component interventions are well defined. An exposure example distinguishes lower risk per trip from fewer expected harms.
\end{abstract}
\section{A declared money-metric choice model}
Fix a baseline population of women, their home locations, travel period, and finite job menus $J_i$. Job $j$ has wage $w_j$, fare $p_{ij}$, time $t_{ij}$, and harm probability $r_{ij}$. Let $v_i\ge0$ value an hour and let $h_i\ge0$ be the expected monetary-equivalent harm conditional on an incident. Assume the expected disutility is $h_ir_{ij}$; nonlinear risk preferences would require a different utility function.

Normalize the nonemployment option to zero and write
\[
 U_i(z)=\max\left\{0,\max_{j\in J_i}
   [w_j-p_{ij}(z)-v_it_{ij}(z)-h_ir_{ij}(z)]\right\}.
 \tag{1}
\]
The same money metric is used across policies for the same individual. A distributional welfare sum additionally needs interpersonal weights. Employment occurs when a job strictly beats zero, with a fixed tie rule assigning nonemployment at equality.
\section{A pure safety intervention on a fixed menu}
Suppose policy one weakly reduces every feasible job's risk and leaves all other attributes unchanged. Let
$\Delta_i=\max_{j\in J_i}h_i[r_{ij}(0)-r_{ij}(1)]$, with zero when the menu is empty.

\begin{proposition}
In the fixed-menu model,
\[
 0\le U_i(1)-U_i(0)\le\Delta_i,\qquad E_i(1)\ge E_i(0).
 \tag{2}
\]
\end{proposition}
\begin{proof}
Every job value weakly rises while the outside option stays fixed, so their maximum weakly rises. Each job value rises by at most $\Delta_i$, so their maximum rises by at most that amount, proving the welfare bound. If the person previously chose employment, some job had strictly positive value. That same job remains strictly positive, so employment continues.
\end{proof}
The proof permits switching jobs; it does not require comparing only the old commute. It does require that the original job remain feasible and that wages and outside options stay fixed. Household bargaining, rationing, and job displacement are absent. Equality in employment can occur even with a strict welfare gain for every initially employed woman, so lack of employment expansion does not prove lack of benefit.

For example, a job paying ten with fare one, valued travel time two, and risk loss three has net value four. Halving the risk loss raises value to $5.5$ without changing employment. A second woman with the same attributes but wage five moves from net value minus one to $0.5$ and enters employment. Both comparisons use fixed illustrative units.
\section{Wage adjustment changes the incidence}
For an equilibrium benchmark, consider a unit mass of workers with outside opportunity cost $\theta$ uniform on $[0,1]$. Every job has the same generalized commuting cost $c$, including valued safety risk, time, and fare under a stated accounting convention. Labor supply is $N=w-c$ in the interior $0<w-c<1$. Employers demand $N=a-bw$ with $b>0$. Market clearing gives
\[
 w(c)=\frac{a+c}{1+b},\qquad
 N(c)=\frac{a-bc}{1+b}.
 \tag{3}
\]
These formulas assume parameters keep supply and demand interior.

A pure safety improvement lowering generalized cost by $\varepsilon>0$ reduces wages by $\varepsilon/(1+b)$ and raises employment by $b\varepsilon/(1+b)$. The net wage available to compensate outside opportunities, $w-c$, rises by the latter amount. Thus the fixed-wage welfare gain overstates the increase in workers' net job surplus.

The equations follow by solving $w-c=a-bw$. Differentiating them establishes the comparative statics. They are a partial-equilibrium labor-market model, not a theorem for all commuting programs. Different labor demand, employer location choices, or bargaining can give other incidence patterns.

For a limiting fixed-job-capacity benchmark, let demand be exactly $N=a\in(0,1)$ regardless of wage. Clearing requires $w=c+a$. Lower commuting cost then lowers wages one for one, while employment and workers' net surplus remain unchanged. Employers gain the wage reduction. If the lower cost is a genuine reduction in expected harm rather than a transfer, aggregate welfare can still rise: the incidence has shifted to employers through equilibrium wages.

For $a=0.4$ and a risk-cost reduction $0.1$, this limiting model keeps employment at $0.4$, lowers wages by $0.1$, and increases employer surplus by $0.04$. Tracking only employment or workers' net wages misses that total gain under equal welfare weights.
\section{A safety contribution needs component interventions}
Let $s\in\{0,1\}$ index a feasible safety-only component and $c\in\{0,1\}$ a feasible time/fare component. Define population outcome $W(s,c)$ under the same equilibrium-selection convention. The combined effect decomposes along one path as
\[
 W(1,1)-W(0,0)
 =[W(1,0)-W(0,0)]+[W(1,1)-W(1,0)].
\]
Reversing the order changes the assigned safety contribution by
\[
 I=W(1,1)-W(1,0)-W(0,1)+W(0,0).
 \tag{4}
\]
Both identities follow by addition and subtraction. If the safety component interacts with fares or time, there is no unique path-free attribution without an additional convention.

A randomized two-by-two market design can identify all four cell means if the interventions are feasible, assignment is positive in every cell, and cross-market spillovers are negligible. If adding security personnel necessarily changes travel time or capacity, the hypothetical safety-only cell may not exist. A regression controlling for realized commuting time does not create it; time is then itself a post-policy choice or mediator.

Average effect estimates should use the same baseline eligible population, including nonworkers and people who move. Conditioning on observed commuters selects a policy-dependent subgroup. The same concern applies to an analysis restricted to women who retain their original jobs.
\section{Risk per trip and total exposure}
Suppose a woman initially makes one relevant trip with incident probability $0.2$. After a safer and more attractive service is introduced she makes three trips, each with probability $0.1$. Expected incident counts rise from $0.2$ to $0.3$ by linearity, even though every trip is safer. Independence is unnecessary for expected counts; it would matter for the probability of at least one incident.

This is not evidence that the safer policy lowers welfare. Additional trips may provide employment or other benefits exceeding their residual risks. It shows why a safety outcome must specify probability per trip, cumulative incident probability, expected severity, and exposure. A lower per-trip rate cannot be substituted for a lower total harm burden without measuring travel behavior.
\section{Remaining research question}
The fixed-menu monotonicity theorem and the wage-incidence model solve separate restricted problems. They do not identify actual risk valuations, labor demand, household constraints, displacement, or the causal safety component of a bundled transport policy. The welfare accounting also needs fares treated as transfers or resource payments consistently with the operator boundary.

The broad catalogue agenda remains unresolved. A credible application must follow baseline women, compare implemented bundles, measure both safety and exposure, and distinguish individual access effects from equilibrium job and wage responses. No clinical, legal, or transport recommendation is inferred from these illustrative models.

\section{Completion Estimate}
45\%

This is a post hoc, heuristic estimate for the full catalogue target.
The attempt proves fixed-menu safety monotonicity, derives wage-incidence and capacity examples, and specifies feasible safety-cost decompositions and cumulative exposure accounting. Full commuting policy effects still require measured harm valuations, baseline-woman follow-up, credible component interventions, and endogenous job, wage, household, and transport responses.

\begin{thebibliography}{9}
\bibitem{flfp} R. Heath and coauthors (2025).
\emph{Female Labour Force Participation}, VoxDevLit 11(2).
\url{https://voxdev.org/voxdevlit/female-labour-force-participation}.
The primary review motivates safety and employment research; the particular commuting and labor-demand formulas are explicit assumptions.
\end{thebibliography}

\end{document}
